\documentclass[10pt,a4paper]{amsart}
\usepackage[margin=19mm]{geometry}
\usepackage{amsmath,amssymb,mathtools}
\usepackage[scr=rsfs,cal=euler]{mathalfa}
\usepackage{xfrac}
\usepackage[unicode,hidelinks,bookmarks=false]{hyperref}
\newcommand{\ie}{i.e.\ }
\newcommand{\cf}{cf.\ }
\newcommand{\resp}{respectively\ }
\newcommand{\Subsection}{\subsection}
\providecommand{\nobreakdash}{\nobreak\hbox{-}}
\setlength{\emergencystretch}{2em}
\multlinegap0pt
\usepackage{mathtools}
\usepackage[all]{xy}
\usepackage{amssymb}
\usepackage{enumerate}
\newtheorem{theorem}{Theorem}
\newtheorem{lemma}{Lemma}
\newcommand{\CMS}{\operatorname{\underline{CM}}}
\newcommand{\D}{\mathcal{D} }
\newcommand{\Dc}{\mathcal{D}}
\newcommand{\Fc}{\mathcal{F} }
\newcommand{\Gc}{\mathcal{G} }
\newcommand{\Hc}{\mathcal{H} }
\newcommand{\Hom}{\operatorname{Hom}}
\newcommand{\Z}{\mathbb{Z} }
\newcommand{\cone}{\operatorname{cone}}
\newcommand{\eR}{\mathcal{R} }
\newcommand{\m}{\mathfrak{m} }
\newcommand{\ov}{\overline}
\newcommand{\rt}{\rightarrow}
\newcommand{\wh}{\widehat }
\title{On the Stable category of maximal Cohen-Macaulay modules over Gorenstein rings-II}
\author{Tony J. Puthenpurakal}
\date{Source: arXiv:2609.18626v1 (CC BY 4.0)}
\begin{document}
\maketitle

\noindent\emph{Source and licence.} On the Stable category of maximal Cohen-Macaulay modules over Gorenstein rings-II. Version arXiv:2609.18626v1 (CC BY 4.0), distributed by arXiv under the Creative Commons Attribution 4.0 International licence, http://creativecommons.org/licenses/by/4.0/, as recorded in the arXiv licence metadata for this exact version and shown on the abstract page https://arxiv.org/abs/2609.18626v1. Full licence notice: \texttt{licenses/arxiv-2609.18626-CC-BY-4.0.txt}.\par\smallskip

\noindent\textit{Editorial scope.} Counted unit: the article's theorem fourth (source0.tex lines 214--217). The article's complete original proof of it is delivered with it (source0.tex lines 930--1000, one \texttt{proof} environment of 71 lines, which the article places away from the statement). No mathematics is written, repaired or re-derived: the statement and the proof are the article's own lines, in the article's own notation, and no retained line is substituted or re-flowed.  Retained uncounted as the source-local context the proof needs: lines 919--922.  Nothing was omitted from any retained span, so the package carries no omission record.  No retained line contains a citation, so the wrapper carries no bibliography and no bibliography entry.  No retained line contains a figure environment, an image-inclusion command or drawing code, so no image is delivered and the package declares no asset dependency.  Editorial formatting only, in addition: an AMS \texttt{amsart} preamble that restates the article's own declarations and macros used by the retained lines, namely the environments \texttt{theorem}, \texttt{lemma} on separate counters, together with the standard packages those lines need.  The retained environments print in this document as Theorem 1, Lemma 1 in order of appearance, whereas the article prints the same items with the numbering of its own full document.  The complete native source of the article as distributed in the arXiv e-print for this exact version is bundled unchanged as \texttt{sources/210710/source0.tex}.
\begin{lemma}\label{clinton}
Let $V = \bigoplus_{n \in \Z}V_n$ be a finitely generated graded $\eR = \eR_A(I)$-module with $V_n = 0$ for all $n \ll 0$.
Let $0 \rt Y \rt X \rt V \rt 0$ be a MCM-approximation of $V$. Then $X = \eR(\Fc, H)$ where $H$ is a free $A$-module and $\Fc$ is an $I$-stable filtration on $H$.
\end{lemma}

\begin{theorem}
\label{fourth} Let $(A,\m)$ be a Gorenstein local ring and let $I \subseteq \m$ be an ideal such that $\eR_A(I)$ is Gorenstein. Then  there exists a triangulated functor \\ $\Psi \colon \CMS_\Z(\eR_A(I)) \rt \CMS(A)$ such that $\Psi$ induces an equivalence
$$\ov{\Psi} \colon \frac{\CMS_\Z(\eR_A(I))}{\ker \Psi} \cong \CMS(A).$$
\end{theorem}

\begin{proof}[Proof of Theorem \ref{fourth}]
As $\Psi$ is dense and full it follows that $\ov{\Psi}$ is dense and full. Thus it suffices to show $\ov{\Psi}$ is faithful. Set $\Dc =  \CMS_\Z(\eR_A(I))/\ker \Psi$. Suppose
$\beta \in \Hom_{\Dc}(\eR(\Fc, M), \eR(\Gc, N)$ be such that $\ov{\Psi}(\beta) = 0$.
Note  $\beta $ can be written as a left fraction
\[
\xymatrix{
\
&\eR(\Hc, E)
\ar@{->}[dl]_{u}
\ar@{->}[dr]^{f}
 \\
\eR(\Fc, M)
\ar@{->}[rr]_{\beta = fu^{-1}}
&\
&\eR(\Gc, N)
}
\]
with $\cone(u) \in \ker \Psi$. It follows that $\Psi(f) = 0$. Thus it suffices to show that
$f = 0$ in $\D$. We note that as $\Psi(f) = 0$ we get that $f_* \colon E \rt N$ factors through a free $A$-module.
After shifting we may assume that $\Hc_n = E$ for $n \leq 0$ and $\Hc_1 \neq E$.

Claim-1: We may assume that $\Gc_n = N$ for $n \leq 0$.

(Note we are not asserting $\Gc_1 \neq N$.)

Proof of Claim-1. Suppose $\Gc_n = N$ for $n \leq a$ with $a < 0$.
We note that $\Gc_n \subseteq \Gc_{n+a}$ for all $n \in \Z$ as $a < 0$. So we have an inclusion of MCM $\eR$-modules
$ \eR(\Gc, N) \xrightarrow{i} \eR(\Gc(a), N)$, say  with co-kernel $V$. We note that $V_n =0$ for $n \ll 0$. Taking  MCM approximations we have a triangle in  $\CMS_\Z(\eR_A(I))$
$$ \eR(\Gc, N) \xrightarrow{i}  \eR(\Gc(a), N) \rt X \rt \Omega^{-1}_{\eR}(\eR(\Gc, N)).$$
We note that by \ref{clinton} it follows that $X \in \ker \Psi$. So $i$ is an isomorphism in $\D$.
Thus it suffices to show that $g = i\circ f = 0$ in $\D$. Notice $g_* = f_* =0$ in $\CMS(A)$.


We have an exact sequence $0  \rt \eR(I, E) \xrightarrow{i_E} \eR(\Hc, E) \rt U \rt 0$. We note that $U_n = 0$ for $n \ll 0$.  Let $Y_E$ be MCM approximation of $\eR(I, E)$. Taking MCM approximations we get by \ref{clinton} that $Y_E \cong \eR(\Hc, E)$ in $\D$. Similarly if $Y_N$ is MCM approximation of $\eR(I, N)$ then $Y_N \cong \eR(\Gc, N)$.

We note that $f_* \colon E \rt N$ induces $\wh{f_*} \colon \eR(I, E) \rt \eR(I, N)$. We have an induced map $X(\wh{f}) \colon Y_E \rt Y_N$. As $f_*$ factors through a free $A$-module say $W$ we get that $\wh{f_*}$ factors through $\eR(I, W)$. So $X(\wh{f_*}) $ factors through $\eR(I, W)$. Thus $X(\wh{f_*}) = 0$ in $\CMS_\Z(\eR_A(I))$.

Note we have a commutative diagram
\[
  \xymatrix
{
 \eR(I,E)
\ar@{->}[r]^{i_E}
\ar@{->}[d]^{\wh{f_*}}
 & \eR(\Hc, E)
\ar@{->}[d]^{f}
\\
 \eR(I, N)
 \ar@{->}[r]^{i_N}
 & \eR(\Gc, N)
\
 }
\]
Taking MCM approximations we get a commutative diagram in $\CMS_\Z(\eR_A(I))$:
\[
  \xymatrix
{
 Y_E
\ar@{->}[r]^{X(i_E)}
\ar@{->}[d]^{X(\wh{f_*})}
 & \eR(\Hc, E)
\ar@{->}[d]^{f}
\\
 Y_N
 \ar@{->}[r]^{X(i_N)}
 & \eR(\Gc, N)
\
 }
\]
As discussed earlier $X(i_E)$ and $X(i_N)$ are isomorphisms in $\D$. Also $X(\wh{f_*}) = 0$ in $\CMS_\Z(\eR_A(I))$ and hence in $\D$. Thus $f = 0$ in $\D$. The result follows.
\end{proof}

\end{document}
